\Needspace{20\baselineskip}\vspace{6mm}
\pdfbookmark[1]{Die neue Einrichtung des Marburger mathematisch-physikalischen Instituts.}{gerling-institut-1842}
\begin{center}\small Göttingische gelehrte Anzeigen.\quad 1842.\quad April 4.\end{center}

Professor \textsc{Gerling} in Marburg, Korrespondent der königlichen Sozietät, hat Hofrat \textsc{Gauss} einen Brief mit einem Bericht über
\begin{center}\textit{die neue Einrichtung des dortigen mathematisch-physikalischen Instituts}\end{center}
geschrieben. Wir geben daraus hier umso lieber einen Auszug, als bei der Einrichtung sowohl auf die gewöhnlichen Bedürfnisse des Unterrichts als auch auf die Wissenschaft selbst Rücksicht genommen wurde und die Räumlichkeiten mehrere Besonderheiten aufweisen.

Das Institut befindet sich seit vergangenem Herbst in einem Staatsgebäude, das auf Befehl Seiner Hoheit des Kurprinzen und Mitregenten umgebaut und der Universität überlassen wurde, dem sogenannten Dörnberger Hof. Das zweite Stockwerk enthält außer dem Hörsaal vier helle Säle, ein kleineres Zimmer, das für optische Versuche verdunkelt werden kann, und noch zwei eigene Arbeitszimmer. Die Instrumente sind mit so viel Platz aufgestellt, dass jedes unmittelbar an seinem Platz benutzt werden kann, ohne den gleichzeitigen Gebrauch der anderen zu verhindern oder zu
\SourceLeaf{603}stören, und dass Studierende mit entsprechender Neigung und ausreichender Vorbereitung bequem eigene Übungen ausführen können. Das Institut besitzt übrigens bereits eine reichhaltige Ausstattung mit physikalischen Geräten, die jedes Jahr weiter vervollständigt und erweitert wird. Bei geöffneten Zwischentüren bietet diese Zimmerreihe eine freie gerade Linie von 150 Fuß Länge, ein in der Tat sehr wertvoller Vorteil für mancherlei Zwecke.

Die Dienstwohnung des Direktors befindet sich im dritten Stockwerk. Das Institut hat auch eine eigene Werkstatt und Schmiede, die im Erdgeschoss untergebracht sind.

Für Versuche, die im Freien angestellt werden müssen, gibt es sowohl in der Umgebung des Gebäudes als auch auf der Plattform eines auf Fels gegründeten und unmittelbar mit dem Gebäude verbundenen Turms hinreichend Gelegenheit. Der Turm bietet zugleich ein für manche physikalische Arbeiten überaus wertvolles Hilfsmittel, nämlich eine freie Fallhöhe von etwa 80 Fuß. Um sie zu erreichen, wurden in den Fußböden der drei Stockwerke des Turms und im Boden über dem Keller (einem ehemaligen Verlies) quadratische Öffnungen angebracht. Auf der oberen Plattform, zu der vom Hausdach aus eine Wendeltreppe führt, befindet sich außerdem ein achteckiger Pavillon von 15 Fuß Durchmesser, in dessen Fußboden eine ähnliche Öffnung ist. die je nach Umständen mit einer Falltür abgedeckt oder mit einer Galerie umgeben werden kann. Im Turm findet sich auch eine sehr feste, beinahe 20 Fuß hohe Mauer, die bei manchen Gelegenheiten wichtige Vorteile bieten kann.

Im Zusammenhang mit dieser Einrichtung wurde auch ein Bedarf gedeckt, der an einer Universität ohne Sternwarte besonders spürbar ist. Dort konnten zum Beispiel die heute für so viele physikalische Arbeiten unentbehrlichen Zeitbestimmungen bisher immer nur durch die zeitraubende Methode korrespondierender Sonnenhöhen gewonnen werden. Dazu brauchte man nur die Türen und Fenster des Pavillons passend einzurichten, um ihn für alle erforderlichen Beobachtungen mit den beweglichen astronomischen Instrumenten nutzbar zu machen, die das Institut zum Teil schon lange besaß. Ein tragbares Passageninstrument von \textsc{Ertel} hat zum Beispiel von seinem gewöhnlichen Standort aus einen völlig freien Bewegungsbereich, im Meridian von Horizont zu Horizont und im ersten Vertikal vom westlichen Horizont bis zu einer östlichen Zenitdistanz von etwa 11 Grad. Die Lage dieses Platzes
\SourceLeaf{604}wurde durch die von Professor \textsc{Gerling} ausgeführte und an die hannoversche Gradmessung angeschlossene trigonometrische Vermessung des Kurfürstentums bestimmt. Ihrer ausführlichen Veröffentlichung sehen wir mit Spannung entgegen. Die Breite beträgt $50^0\;48'\;46''9$, die Länge von Ferro $26^0\;26'\;2''3$.

Schließlich ist noch eine Einrichtung zu erwähnen, die als wesentlich für die Vollendung des ganzen Plans galt. So ausgezeichnet sich das Hauptgebäude auch für alle übrigen Zwecke eignete, hatte es doch den Nachteil, dass es sich verhältnismäßig am wenigsten für meteorologische Beobachtungen nutzen ließ, da es vom erheblich höheren Schlossberg überragt wird. Außerdem ist es für mancherlei Zwecke sehr wünschenswert, zwei ziemlich weit voneinander entfernte und doch gegenseitig erreichbare Räumlichkeiten zur Verfügung zu haben. Diesem doppelten Bedürfnis half man ab, indem ein vorhandenes altes Türmchen auf dem höchsten Punkt des Schlossbergs, in gerader Linie 1900 Fuß entfernt und etwa 100 Fuß höher gelegen, zu einem meteorologischen Turm ausgebaut wurde. So kann unter anderem abseits der Stadt und damit unter Ausschluss jeder denkbaren Gefahr ein Blitzableiter für Beobachtungen eingerichtet werden. Damit ist auch die Möglichkeit gegeben, dort demnächst zum Beispiel Versuche zur magnetischen Telegraphie anzustellen.

